% Generated draft: requires mathematical and source review.
\documentclass[11pt]{article}
\usepackage[margin=25mm]{geometry}
\usepackage{fontspec}
\setmainfont{Latin Modern Roman}
\usepackage{amsmath,amssymb,mathtools,unicode-math}
\setmathfont{Latin Modern Math}
\usepackage{enumitem,needspace,xurl,hyperref}
\setlength{\parindent}{0pt}
\setlength{\parskip}{5pt}
\setlength{\emergencystretch}{3em}
\newcommand{\N}{\mathbb{N}}
\newcommand{\Z}{\mathbb{Z}}
\newcommand{\Q}{\mathbb{Q}}
\newcommand{\R}{\mathbb{R}}
\newcommand{\C}{\mathbb{C}}
\newcommand{\F}{\mathbb{F}}
\newcommand{\A}{\mathbb{A}}
\newcommand{\PP}{\mathbb{P}}
\newcommand{\E}{\mathbb{E}}
\newcommand{\eps}{\varepsilon}
\newcommand{\dd}{\,\mathrm{d}}
\begin{document}
\section*{UnsolvedMath ID: 600001}
\textbf{Source problem code:} catalogue entry
\subsection{Definitions and mathematical statement}
For a convex domain with a sufficiently regular boundary, its inner billiard map sends an oriented incoming line to the outgoing oriented line obtained by specular reflection at the next boundary impact; this definition is restricted to impacts where the reflection and the choice of next impact are well defined. Two planar ellipses are confocal when they have the same pair of foci. For planar outer billiards, fix one of the two consistent tangent orientations: if the selected tangent from an exterior point $x$ touches the boundary at $t$, the map is $x\mapsto 2t-x$. Concentric homothetic ellipses have a common center and differ by a positive dilation about that center.

\textbf{Planar inner billiards (the source conjecture).} Let $D_1\subset D_2\subset\mathbb R^2$ be nested convex domains, and let $T_1,T_2$ denote their inner billiard maps on oriented lines intersecting both domains, wherever the maps and their compositions are defined. Does
\[
T_1\circ T_2=T_2\circ T_1
\]
imply that $\partial D_1$ and $\partial D_2$ are confocal ellipses? The source also asks for the corresponding questions for inner and outer billiards in higher dimensions.

\textbf{Branch-specific information supplied with the record.} The research summary reports the planar inner conclusion for nested convex domains with piecewise $C^4$ boundaries (Glutsyuk, 2017), and the inner conclusion that the boundaries are confocal ellipsoids for nested strictly convex $C^2$ billiards in dimension $d\geq3$ (Glutsyuk, 2020). It attributes the latter approach to Berger's theorem on quadrics admitting caustics in dimension at least three, and reports extensions to constant-curvature spaces. These are reports of results under the stated regularity hypotheses, not a strengthening to all convex domains. For planar outer billiards the reported conclusion from commutation is instead that the boundaries are concentric homothetic ellipses (Tabachnikov, 1994). The higher-dimensional outer question has no resolution identified in the supplied summary and is described there as apparently still open; its precise billiard model is not specified.
\subsection{Short English statement}
Ask whether commutation of the inner billiard maps for nested planar convex domains forces confocal ellipses, and formulate the corresponding multidimensional inner and outer questions. The supplied summary reports affirmative inner results under stated regularity hypotheses, a distinct planar outer conclusion of concentric homothetic ellipses, and no identified resolution of the higher-dimensional outer branch.
\subsection{Sources}
\begin{itemize}
\item S. Tabachnikov, \emph{A Baker's Dozen of Problems} (2015), Conjecture 1, source list and publisher HTML: \url{https://amj.math.stonybrook.edu/html-articles/Files-2015-2024/14-01/}; accessed 2026-07-29.
\item Supplied research-summary attributions: Glutsyuk (2017), planar inner commuting billiards; Glutsyuk (2020), higher-dimensional inner commuting billiards and constant-curvature extensions; Tabachnikov (1994), planar commuting outer billiards; Berger's theorem on caustics and quadrics in dimensions at least three. The record supplies no further bibliographic identifiers for these attributions.
\item Local source record: \texttt{data/unsolvedmath/problems.json}, permanent ID $600001$, problem number catalogue entry, \emph{Baker's Dozen --- Commuting billiard maps}. The record labels the problem solved and SOLVED-IN-LITERATURE, while its summary separately leaves the higher-dimensional outer branch apparently open; its background marks source-status and rights checks NEEDS\_REVIEW.
\end{itemize}
\end{document}
